\documentclass{article}
\usepackage{fullpage}

%load needed packages
\usepackage{graphicx}
\usepackage{array}
\usepackage{booktabs}
\usepackage[utf8]{inputenc}
\usepackage[T1]{fontenc}

\usepackage[spanish]{babel} % Paquete para el idioma español
\usepackage{float}  % Necesario para [H]
\usepackage{listings}
\usepackage{xcolor}

\definecolor{codegreen}{HTML}{5AB2FF}
\definecolor{morado}{HTML}{AD88C6}
\definecolor{BG}{HTML}{EEEEEE}
\definecolor{azul}{HTML}{4D869C}
\definecolor{sqlblue}{HTML}{FF8C00} % Color para las palabras clave SQL

% Estilo para DDL
\lstdefinestyle{ddlstyle}{
	language=SQL,
	backgroundcolor=\color{BG},
	commentstyle=\color{codegreen},
	basicstyle=\ttfamily\small,
	keywordstyle=\color{azul},
	stringstyle=\color{morado},
	showstringspaces=false,
	breaklines=true,
	frame=shadowbox,
	numbers=left,
	numberstyle=\tiny\color{gray},
	captionpos=b,
}

% Estilo para SQL
\lstdefinestyle{sqlstyle}{
	language=SQL,
	backgroundcolor=\color{BG},
	commentstyle=\color{codegreen},
	basicstyle=\ttfamily\small,
	keywordstyle=\color{sqlblue}, % Color diferente para palabras clave SQL
	stringstyle=\color{morado},
	showstringspaces=false,
	breaklines=true,
	frame=shadowbox,
	numbers=left,
	numberstyle=\tiny\color{gray},
	captionpos=b,
}

\begin{document}



% Portada
\begin{titlepage}
	\centering
	\vspace*{3cm}
	
	% Título destacado
	{\Huge \textbf{Diseño y Explotación de un almacén de 
			UCI Sanitaria}\\[0.5cm]}
	
	% Espacio y logotipo (si lo tienes, por ejemplo el logo de tu universidad)
	\vspace{2cm}
	\includegraphics[width=0.3\textwidth]{images/uma_logo.jpg}\\[1cm]
	
	% Nombre del autor
	{\LARGE \textbf{Alejandro Silva Rodríguez}\\[0.5cm]}
	{\LARGE \textbf{Marta Cuevas Rodríguez}\\[0.5cm]}
	{\large \textit{Almacenes De Datos}\\
		Universidad de Málaga\\
		}
	
	\vfill
	
	% Fecha en la parte inferior de la página
	{\large Septiembre 2024}
\end{titlepage}

% indice
\tableofcontents

\newpage
\section{Introducción}
\label{sec:introduccion}
Breve introducción sobre el objetivo del proyecto y el contexto en la UCI Sanitaria. Descripción de la importancia de un almacén de datos para el análisis de datos de pacientes, su relevancia en la asignatura de Almacenes de Datos, y el tipo de información que se modelará (por ejemplo, pacientes con enfermedades cardiovasculares).

\section{Diseño Conceptual}
\label{sec:diseno_conceptual}
\subsection{Descripción del Modelo Conceptual}
Presenta el diseño conceptual del almacén de datos. Explica las entidades y dimensiones seleccionadas, y proporciona un diagrama conceptual generado en la herramienta elegida. Justifica las decisiones de diseño, destacando aspectos como la selección de dimensiones y las jerarquías establecidas.

\subsection{Decisiones de Diseño}
Explica y justifica las decisiones tomadas durante el diseño conceptual, incluyendo cualquier simplificación, agrupamiento o denormalización de las dimensiones. Discute el propósito de cada decisión y cómo contribuye a cumplir con los objetivos del proyecto.

\section{Diseño Lógico}
\label{sec:diseno_logico}
\subsection{Transformación a Modelo en Copo de Nieve / Constelación}
Describe cómo el diseño conceptual se convierte en un modelo lógico de tipo copo de nieve o constelación. Incluye un diagrama del diseño lógico, destacando las relaciones entre las tablas y el esquema de la dimensión tiempo denormalizada.

\subsection{Definición de Claves}
Define las claves primarias y las claves externas de cada tabla en el esquema lógico, explicando cómo cada clave se selecciona y se asocia a las entidades correspondientes.

\section{Creación de Tablas en SQL Server}
\label{sec:creacion_tablas}
\subsection{Generación del Código DDL}
Presenta el código DDL para la creación de las tablas del almacén en SQL Server. Explica las instrucciones y pasos necesarios para ejecutar este código correctamente, incluyendo cualquier configuración especial o instrucción adicional necesaria.

\subsection{Instrucciones para Crear el Almacén}
Proporciona una guía para desplegar el almacén de datos en SQL Server usando el código DDL, detallando los pasos a seguir y cualquier configuración de permisos o ajustes adicionales.

\section{Dificultades Encontradas}
\label{sec:dificultades_encontradas}
Enumera y describe las dificultades enfrentadas durante el desarrollo del proyecto. Explica cómo se resolvieron estos problemas o, en caso de no haber encontrado complicaciones, indícalo claramente.

\section{Conclusión}
\label{sec:conclusion}
Resumen de los principales resultados del proyecto, con un breve comentario sobre la utilidad del almacén de datos desarrollado y cómo este se alinea con los objetivos de la asignatura.


\newpage
\section{Acceso al Repositorio}

Toda la información adicional, incluyendo el código fuente y la documentación completa de este proyecto, está disponible en el repositorio de GitHub \cite{silva2024github}.

% Incluir la bibliografía
\bibliographystyle{plain}  % Estilo de la bibliografía (por ejemplo, plain, alpha, ieee, etc.)
\bibliography{bibli}  % Nombre del archivo .bib sin la extensión

\end{document}
